\chapter{Discusión general y conclusiones}
\label{chap:discusion}

\section{Síntesis}

Esta tesis retoma la pregunta y la hipótesis con las que se aprobó su tema en diciembre de 2025 (Informe de Taller de Tesis I): si el mercado penaliza de forma no lineal la coexistencia de fragilidad bancaria y vulnerabilidad macroeconómica. Esa hipótesis se validó entonces sobre una muestra piloto de cuatro países y una especificación simple, en la que la interacción $JLoss\times GaR$ emergía significativa solo bajo efectos fijos de país y tiempo, y no bajo una regresión \textit{pooled}. El Capítulo~\ref{chap:paper2} pone la misma hipótesis a prueba a una escala mucho mayor —trece economías emergentes, con la fragilidad bancaria construida por un motor de datos propio desde la terminal Bloomberg— y con una batería de identificación considerablemente más exigente que la disponible en la etapa piloto: veinticuatro regresiones anidadas, un modelo de umbral, proyecciones locales, variables instrumentales por \textit{shift-share}, \textit{wild cluster bootstrap} y un test de falsación directo por régimen de crisis.

La respuesta que esa evidencia entrega es más precisa —y por eso más interesante— que la de la muestra piloto. Los canales de nivel se sostienen con solidez: la fragilidad bancaria sistémica y el riesgo de cola del crecimiento elevan cada uno el spread soberano en las doce especificaciones de la batería principal. La complementariedad entre ambos —el coeficiente $\hat\beta_3$ de $JLoss\times D$, con $D\equiv-GaR$— tiene el signo predicho por el informe original y es corroborada por un modelo de umbral, pero no es una regularidad uniforme del panel completo de trece países ($\hat\beta_3=+0{,}16$, $p=0{,}26$). Se identifica, en cambio, en dos cortes que apuntan al mismo mecanismo económico: el núcleo de economías emergentes de financiamiento externo ($\hat\beta_3=+0{,}47$, $p=0{,}023$) y los trimestres sin crisis financiera aguda ($\hat\beta_3\approx+1{,}2$, $p<0{,}01$); combinados, $\hat\beta_3=+0{,}94$ ($p=0{,}004$). Un test de falsación directo —interactuar la complementariedad con un vector de crisis en vez de excluir trimestres— confirma que el mecanismo se apaga específicamente bajo respaldo oficial masivo (crisis financiera global y pandemia) y sobrevive en el estrés emergente de 2015--2016, que no tuvo ese respaldo. La no significancia del panel completo es, en definitiva, un promedio de regímenes distintos: exactamente el tipo de matiz que una muestra piloto de cuatro países y una especificación simple no podía revelar. El Capítulo~\ref{chap:paper2} desarrolla esta evidencia en detalle, incluida la caracterización completa de dónde y cuándo la complementariedad aparece en los precios soberanos, y no se repite aquí.

\section{Implicancias de política}

La lectura de política que emerge de esta investigación es precisa porque delimita cuándo aplica. La regulación orientada a la resiliencia del sistema bancario —capital, liquidez, provisiones contracíclicas— reduce la exposición del soberano a los escenarios macroeconómicos adversos y su prima de riesgo en los mercados internacionales; la fragilidad bancaria y el riesgo de cola del crecimiento deben monitorearse conjuntamente, porque es su coincidencia —y no cada uno por separado— la señal de alarma relevante para el costo de financiamiento soberano. La evidencia de este panel establece esa implicancia con más fuerza para las economías más expuestas al financiamiento externo de cartera, donde la complementariedad es estadísticamente significativa, y para los períodos en que el soberano no cuenta con respaldo oficial explícito o implícito. En las economías con mercados de deuda local profundos, o durante episodios en que operan líneas de \textit{swap} de bancos centrales y financiamiento multilateral de emergencia, la misma fragilidad bancaria se traduce con menor intensidad en un mayor costo de financiamiento soberano —lo que no exime de la necesidad de vigilancia prudencial, pero sí matiza la urgencia relativa de la señal de mercado como gatillo de alarma temprana.

\section{Limitaciones}

El Capítulo~\ref{chap:paper2} (Sección~7.3) detalla ocho limitaciones específicas de la evidencia empírica —entre ellas la falta de significancia sobre el panel completo, la escasez de episodios de cola independientes, la compresión transversal de $JLoss$ por homogeneidad de fuente, la cobertura desigual del EMBI, el carácter estimado de los regresores y los límites de la identificación por variables instrumentales— y no se repiten aquí. A nivel de tesis conviene sintetizar dos de ellas. Primera, la misma compresión transversal de $JLoss$ que limita la identificación de la interacción principal limita, con más fuerza todavía, cualquier interacción de orden superior: la pregunta de si la concentración bancaria amplifica la complementariedad fragilidad--crecimiento —motivada por la literatura de competencia y fragilidad bancaria, que no tiene una predicción unívoca sobre su signo— se explora como agenda futura en el Capítulo~\ref{chap:paper2} (Sección~7.4) precisamente porque el panel actual no tiene potencia estadística para dirimirla con ninguna de las tres proxies de concentración probadas. Segunda, ninguna especificación de este trabajo cierra por completo el problema de simultaneidad del \textit{doom loop} —causalidad inversa desde el spread soberano hacia la fragilidad bancaria vía el valor de la cartera de deuda pública que los bancos mantienen—: la estrategia de variables instrumentales explorada no cierra la identificación sobre esta muestra reducida bajo ninguno de los tres instrumentos probados, y la evidencia del canal de nivel descansa, en consecuencia, en el MCO con efectos fijos y en las proyecciones locales.

\section{Agenda de investigación futura}

La agenda que se desprende de esta investigación tiene dos prioridades, ambas de datos antes que de método. Primero, ampliar el grupo de contraste de la heterogeneidad transversal —hoy solo Polonia e India representan a las economías de deuda local profunda— con más economías de ese perfil, y construir una medida directa de participación extranjera en la deuda soberana que sirva de moderador continuo en lugar de la dicotomía por tipo de economía. Segundo, ampliar el corte transversal de $JLoss$ incorporando a Argentina, Egipto y Rusia —hoy solo les falta el $GaR$ o el EMBI, no la fragilidad bancaria—, lo que además daría potencia por primera vez a la prueba de si la concentración bancaria amplifica la complementariedad, hoy no identificada con ninguno de los tres proxies de concentración disponibles (Capítulo~\ref{chap:paper2}, Sección~7.4). En un segundo orden de prioridad, instrumentos adicionales para el canal de nivel —choques de política monetaria identificados, u otro instrumento cuya relevancia sea específicamente bancaria y no macroeconómica agregada— permitirían cerrar la identificación causal que las variables instrumentales de \textit{shift-share} probadas en este trabajo no lograron cerrar.

\section{Conclusión}

Esta tesis sostiene, con la evidencia y el rigor disponibles hoy, que la fragilidad bancaria sistémica y el riesgo a la baja del crecimiento elevan cada uno el riesgo soberano de las economías emergentes, y que su coincidencia lo amplifica por encima de la suma de sus efectos individuales cuando el mercado espera que el pasivo contingente bancario recaiga sobre un soberano no respaldado. Este hallazgo confirma, con un panel bastante más amplio y una batería de identificación considerablemente más exigente, la hipótesis central con la que se aprobó el tema de esta tesis en el Informe de Taller de Tesis I: que el mercado penaliza de forma no lineal la coexistencia de fragilidad bancaria y vulnerabilidad macroeconómica. Lo hace, sin embargo, con una precisión que la etapa piloto no podía anticipar: esa penalización no es un rasgo uniforme de todas las economías emergentes ni de todos los regímenes de mercado, sino una propiedad de a quién y en qué circunstancias se le fija el precio a la deuda soberana. La evidencia organiza esa propiedad en una batería de veinticuatro regresiones sobre un único panel de trece economías emergentes con la fragilidad bancaria construida desde Bloomberg y el spread medido por el EMBI Global Diversified: establece con solidez los canales de nivel —significativos en las doce especificaciones— y muestra que la complementariedad —del signo predicho y corroborada por un modelo de umbral— se identifica en el núcleo de economías de financiamiento externo ($\hat\beta_3=+0{,}47$) y en los trimestres sin crisis financiera aguda ($\hat\beta_3\approx+1{,}2$); combinados, $+0{,}94$ ($p=0{,}004$), y sobrevive a un test de falsación directo frente a los episodios de respaldo oficial masivo.

Lo que queda establecido son los canales de nivel y el carácter condicional —pero estructural y económicamente interpretable— de la complementariedad entre fragilidad bancaria y riesgo de cola del crecimiento; lo que queda abierto es su generalidad al conjunto de mercados emergentes y su relación con la estructura de competencia del sistema bancario, hoy una pregunta sin potencia estadística para responderse y no un hallazgo. El aporte de esta investigación no es un coeficiente robusto a toda prueba —esos abundan, mal identificados, en la literatura de spreads soberanos—, sino la delimitación precisa de qué puede afirmarse y qué no: la complementariedad fragilidad--crecimiento se identifica en las economías más expuestas al financiamiento externo y en los regímenes sin respaldo oficial, tiene un mecanismo económico que explica por qué se apaga en los demás, y descansa sobre una métrica de fragilidad medida de forma homogénea y reproducible, construida con el motor propio que el Informe de Taller de Tesis I comprometía como trabajo futuro. Esa disciplina —decir dónde el mecanismo muerde y dónde el dato no alcanza— es una contribución a la credibilidad de la evidencia, no solo a su volumen.
